\begin{longtable}{|p{3.2cm}|p{6cm}|p{5.8cm}|}
\hline
Categoria & Descrição sucinta & Estudos representativos \\ \hline
\endfirsthead
\hline
Categoria & Descrição sucinta & Estudos representativos \\ \hline
\endhead
\hline
\endfoot
\endlastfoot
Perfil demográfico e faixa etária &
Idade, gênero, escolaridade e ocupação, com predominância de estudos com idosos e alguns com comparações entre faixas etárias. &
\cite{zaphiris2003age}; \cite{patsoule2012redesigning}; \cite{hart2008evaluating}; \cite{ko2019old}; \cite{xie2010usability}; \cite{ahmad2013teaching}; \cite{lee2021factors}; \cite{sharit2008information}; \cite{romano2013age}; \cite{crabb2014age}; \cite{arfaa2015usability}; \cite{finn2013usability} \\ \hline
Experiência, familiaridade e alfabetização digital &
Experiência com computadores e Web, frequência de uso, distinção entre novatos e experientes e diferentes níveis de literacia digital. &
\cite{cowan2011exploring}; \cite{sharit2008information}; \cite{mastarone2007access}; \cite{gautam2019adding}; \cite{kusumaningayu2017web}; \cite{nosseir2011novice}; \cite{filgueiras2009recoverability}; \cite{delara2016study}; \cite{lum2016development}; \cite{morash2015guiding}; \cite{esticky2023} \\ \hline
Características sensoriais, motoras e cognitivas &
Declínio visual, auditivo, motor e cognitivo, memória, atenção, ansiedade e dificuldades relacionadas à complexidade das tarefas. &
\cite{park2015personalized}; \cite{biswas2014interface}; \cite{wang2020research}; \cite{breaking_barriers}; \cite{taylor2009using}; \cite{wickramathilaka2023addressing}; \cite{sharit2008information}; \cite{crabb2014age}; \cite{romano2013age}; \cite{duarte2025designing}; \cite{esticky2023} \\ \hline
Deficiências e necessidades de acessibilidade &
Deficiências visuais, auditivas, motoras e cognitivas, uso de tecnologias assistivas e necessidade de entrada alternativa ou reorganização do conteúdo. &
\cite{kouroupetroglou2016methodological}; \cite{breaking_barriers}; \cite{halbach2009framework}; \cite{kusumaningayu2017web}; \cite{duarte2025designing}; \cite{gonzalez2025augmenting}; \cite{mowar2025codea11y}; \cite{ramineni2024enhancing} \\ \hline
Características socioculturais &
Contexto cultural, socioeconômico, idioma e participação de cuidadores, com discussão limitada de comparações entre culturas. &
\cite{ko2019old}; \cite{breaking_barriers}; \cite{herrera2026applicability}; \cite{nosseir2011novice}; \cite{dombrovskaia2021guideline}; \cite{lee2021factors}; \cite{ahmad2013teaching} \\ \hline
Contextos de uso, domínio e dispositivo &
Saúde, governo, banking, turismo, educação, mídia social, uso doméstico e dispositivos móveis, smartwatches e computadores. &
\cite{health_problem_solving}; \cite{trilar2021family}; \cite{gatsou2017seniors}; \cite{martinez2020effectiveness}; \cite{duarte2025designing}; \cite{balakrishnan2012user}; \cite{kymäläinen2012creating}; \cite{gonzalez2025augmenting}; \cite{finn2013usability}; \cite{arfaa2015usability}; \cite{lum2016development}; \cite{secure_messaging}; \cite{esticky2023}; \cite{chittaro2003guiding}; \cite{kanjanakuha2019comprehensibility}; \cite{becker2004study} \\ \hline
\caption{Características dos usuários e contextos de uso identificados nos estudos que respondem diretamente à RQ4.}
\label{tab:caracteristicasRQ4}
\end{longtable}